% ============================================
% Assistant Professor Cover Letter – Westmont College
% ============================================

\documentclass[11pt,letterpaper]{letter}

\usepackage{geometry}
\usepackage{microtype}
\usepackage[hidelinks]{hyperref}

\geometry{left=25mm,right=25mm,top=25mm,bottom=25mm}

\signature{Decory Edwards}
\address{
Department of Economics \\
Johns Hopkins University \\
Baltimore, MD 21218 \\
\texttt{dedwar65@jh.edu}
}

\begin{document}

\begin{letter}{
Search Committee \\
Department of Economics and Business \\
Westmont College \\
Santa Barbara, CA
}

\opening{Dear Members of the Search Committee,}

I am writing to express my strong interest in the one-year Assistant Professor position in the Department of Economics and Business at Westmont College, beginning in August 2026. I am a Ph.D.\ candidate in Economics at Johns Hopkins University, advised by Professors Christopher D.~Carroll, Nicholas W.~Papageorge, and Stelios Fourakis, and I expect to receive my degree in May~2026. My research fields are macroeconomics, wealth and income dynamics, and computational economics.

My research agenda focuses on understanding the determinants of wealth inequality and the mechanisms through which heterogeneity in returns to wealth shapes aggregate outcomes. In my job-market paper, I estimate the degree of heterogeneity in individual portfolio returns using micro-data from the Survey of Consumer Finances and simulate a structural model in which households face idiosyncratic return risk. I show that allowing for persistent heterogeneity in returns substantially improves the model’s ability to match the empirical distribution of wealth, offering a tractable way to connect micro-level return dispersion to macroeconomic inequality.

In a second project, I use the Health and Retirement Study to examine the relationship between interpersonal trust and portfolio performance. Building on the literature that finds a hump-shaped relationship between trust and income, I show that a similar relationship holds for individual returns to wealth measured in the HRS data, highlighting a behavioral channel through which social attitudes shape saving and investment outcomes. This work also provides new estimates of the persistent component of returns to net worth, which motivate the calibration of heterogeneity in my structural model.

Westmont College is an especially appealing environment for my teaching and scholarly interests. The department’s emphasis on close faculty student engagement, undergraduate research mentorship, and broad exposure to economic ideas aligns closely with my approach to teaching economics in a liberal arts setting. I would be well prepared to teach courses in macroeconomics and financial markets, and I would welcome the opportunity to contribute to international economics or historically grounded courses that connect economic ideas to broader intellectual traditions.

I am deeply interested in the teacher scholar model and in mentoring undergraduate research. My teaching emphasizes clarity, economic intuition, and the integration of data and real world examples to help students understand both theory and application. In a small classroom environment like Westmont’s, I would aim to foster careful reasoning, respectful discussion, and sustained intellectual curiosity, while supporting students as they develop analytical and writing skills.

I am also drawn to Westmont College’s Christian liberal arts mission and its commitment to educating students within a Protestant evangelical tradition. I value an academic environment that emphasizes moral formation, service, and the integration of faith and learning, and I would be enthusiastic about contributing to a community that seeks to educate students as thoughtful scholars and responsible citizens.

Thank you very much for your time and consideration. I would welcome the opportunity to discuss my application further.

\closing{Sincerely,}

\end{letter}
\end{document}
